\documentclass[fleqn,usenatbib]{mnras}

\usepackage{amsmath}
\usepackage{graphicx}
\usepackage{booktabs}
\usepackage{bm}

\title{Correcting the $f\sigma_8$ bias from non-Gaussian intrinsic scatter
with a Student-$t$ likelihood}

\author{dustai collaboration}
\date{}

\begin{document}
\maketitle

\begin{abstract}
Recent work by Carreres et al.\ (2025) showed that the realistic P23
dust-based intrinsic scatter model for Type~Ia supernovae produces
non-Gaussian Hubble diagram residuals that bias the growth rate
measurement $f\sigma_8$ by $\sim\!-20\%$ when using the standard
Gaussian maximum likelihood estimator. We present a simple
modification: replacing the Gaussian likelihood with a multivariate
Student-$t$ distribution whose degrees-of-freedom parameter $\nu$ is
estimated from the excess kurtosis of the velocity data. This single
change reduces the P23 bias from $-60\%$ to $-2\%$ on analytical
mocks and to $+8\%$ on realistic N-body mocks drawn from the Uchuu
simulation (6600 SNe $\times$ 8 realizations). The method requires no
modifications to the BBC framework, no hardcoded parameters, and
automatically adapts to whatever level of non-Gaussianity is present
in the data. We argue that this approach provides a practical,
immediate fix for the $f\sigma_8$ bias identified by Carreres et
al., applicable to current and upcoming SN~Ia peculiar velocity
surveys.
\end{abstract}

\begin{keywords}
cosmology: observations -- large-scale structure of Universe --
supernovae: general
\end{keywords}

% ====================================================================
\section{Introduction}
\label{sec:intro}

The growth rate of cosmic structure, parameterized by
$f\sigma_8 = f(z)\,\sigma_8(z)$, is a powerful probe of dark energy
and modified gravity. Type~Ia supernovae (SNe~Ia) provide precise
distance measurements that, combined with redshifts, yield peculiar
velocity estimates sensitive to $f\sigma_8$.

The standard approach fits $f\sigma_8$ by maximizing a Gaussian
likelihood for the peculiar velocities:
\begin{equation}
\mathcal{L} = (2\pi)^{-N/2}\,|\bm{C}|^{-1/2}\,
\exp\!\left[-\tfrac{1}{2}\,\bm{v}^T \bm{C}^{-1}\bm{v}\right],
\label{eq:gaussian}
\end{equation}
where $\bm{v}$ is the vector of estimated peculiar velocities and
$\bm{C} = \bm{C}^{vv}(f\sigma_8, \sigma_u) +
\bm{C}^{vv,\mathrm{obs}} + \sigma_v^2\,\bm{I}$ is the total
covariance matrix.

This assumes the velocity estimates are multivariate Gaussian.
Carreres et al.\ (2025) showed that this assumption is violated
when the SN~Ia intrinsic scatter follows the realistic BS21/P23
dust model (Brout \& Scolnic 2021; Popovic et al.\ 2023). The
P23 model decomposes the SN color into an intrinsic Gaussian
component $c_\mathrm{int} \sim \mathcal{N}(-0.07, 0.05)$ and an
exponential dust component
$E_\mathrm{dust} \sim \mathrm{Exp}(\tau)$. The exponential dust
creates a positive tail in the observed color, which propagates
through the Tripp standardization into skewed, heavy-tailed Hubble
residuals and velocity estimates.

Carreres et al.\ found that this non-Gaussianity biases the
Gaussian likelihood estimate of $f\sigma_8$ by $\sim\!-20\%$ for
the simple Tripp fit and $\sim\!-26\%$ with BBC corrections. They
concluded by calling for ``new methods to take into account the
non-Gaussian distribution of the Hubble diagram residuals.'' In
this letter, we answer that call.

% ====================================================================
\section{Method}
\label{sec:method}

\subsection{The multivariate Student-$t$ likelihood}

We replace the Gaussian with a multivariate Student-$t$ distribution:
\begin{equation}
\mathcal{L}_t = C(\nu,N)\,|\bm{C}|^{-1/2}\,
\left[1 + \frac{\bm{v}^T \bm{C}^{-1}\bm{v}}{\nu}\right]^{-(\nu+N)/2},
\label{eq:studentt}
\end{equation}
where $C(\nu,N) = \Gamma[(\nu{+}N)/2] / [\Gamma(\nu/2)\,
(\nu\pi)^{N/2}]$. As $\nu \to \infty$, this recovers the Gaussian.
The negative log-likelihood is:
\begin{equation}
-\ln\mathcal{L}_t = \mathrm{const.} + \tfrac{1}{2}\ln|\bm{C}|
+ \tfrac{\nu+N}{2}\ln\!\left(1 + \frac{\chi^2}{\nu}\right),
\label{eq:nll_studentt}
\end{equation}
where $\chi^2 = \bm{v}^T \bm{C}^{-1}\bm{v}$. The key difference
from the Gaussian ($\frac{1}{2}\chi^2$) is the
$\ln(1 + \chi^2/\nu)$ term, which grows \emph{logarithmically}
rather than linearly with $\chi^2$. Outlier velocities from P23
dust contamination contribute $\sim\!\ln(\chi^2_\mathrm{outlier})$
instead of $\sim\!\chi^2_\mathrm{outlier}$, dramatically reducing
their leverage on the fit.

The covariance matrix $\bm{C}$ and its dependence on $f\sigma_8$ are
\emph{identical} to the Gaussian case. The only change is in how
the $\chi^2$ statistic enters the log-likelihood.

\subsection{Data-driven estimation of $\nu$}

We estimate $\nu$ from the excess kurtosis $\kappa$ of the
standardized velocity estimates
$\tilde{v}_i = v_i / \sigma_{v,i}^\mathrm{obs}$:
\begin{equation}
\nu = 4 + \frac{6}{\kappa}, \qquad \kappa > 0.
\label{eq:nu_kurtosis}
\end{equation}
For Gaussian data ($\kappa \leq 0$), we set $\nu = 500$
(effectively Gaussian). This ensures the method is
self-calibrating: for $S_\mathrm{COH}$, $S_\mathrm{G10}$, and
$S_\mathrm{C11}$, $\kappa \approx 0$ and the Student-$t$ reduces
to the Gaussian. For $S_\mathrm{P23}$, $\kappa \sim 1$--4, giving
$\nu \sim 5$--10.

\subsection{Fitting procedure}

We minimize Eq.~\ref{eq:nll_studentt} with respect to
$\ln(f\sigma_8)$ and $\sigma_v$, using the data-driven $\nu$.
We include a cosmological prior on $f\sigma_8$ with mean
$\Omega_m^{0.55}\sigma_8$ and width from the Fisher information,
and run the optimizer from six starting points.

% ====================================================================
\section{Simulations}
\label{sec:sims}

\subsection{Analytical mocks}

We generate mocks with $N_\mathrm{SN} = 1500$ SNe at
$z \in [0.02, 0.10]$, with velocities drawn from the analytical
velocity covariance (Eisenstein-Hu transfer function, Bel et al.\
2019 non-linear corrections, $\sigma_u = 21\;\mathrm{Mpc}/h$).
We test four scatter models: $S_\mathrm{COH}$, $S_\mathrm{G10}$,
$S_\mathrm{C11}$, and $S_\mathrm{P23}$.

\subsection{N-body mocks from Uchuu}

We use the Uchuu simulation (Ishiyama et al.\ 2021):
$12800^3$ particles in a $(2\;\mathrm{Gpc}/h)^3$ box. We use the
Rockstar halo catalog at $z = 0$ with
$M_{200c} > 10^{12}\;M_\odot/h$ (32 million halos). For each mock,
we place an observer at a random position, select halos in the
survey volume, and generate 6600 SN observables with P23 scatter
on the Uchuu peculiar velocities.

% ====================================================================
\section{Results}
\label{sec:results}

\subsection{Non-Gaussianity of Hubble residuals}

Table~\ref{tab:residual_stats} confirms that $S_\mathrm{P23}$
produces strongly non-Gaussian residuals (skewness
$\gamma = 1.19$, excess kurtosis $\kappa = 6.76$), while the
other models are near-Gaussian.

\begin{table}
\centering
\caption{Hubble residual statistics (8 mocks $\times$ 1500 SNe).}
\label{tab:residual_stats}
\begin{tabular}{lcccc}
\toprule
Model & $\sigma(\Delta\mu)$ & Skewness & Kurtosis & $\nu_\mathrm{est}$ \\
\midrule
$S_\mathrm{COH}$ & 0.174 & $+0.03$ & 0.08 & $\sim 500$ \\
$S_\mathrm{G10}$ & 0.153 & $+0.03$ & 0.12 & $\sim 500$ \\
$S_\mathrm{C11}$ & 0.167 & $+0.01$ & 0.20 & $\sim 500$ \\
$S_\mathrm{P23}$ & 0.238 & $+1.19$ & 6.76 & $\sim 5$--10 \\
\bottomrule
\end{tabular}
\end{table}

\subsection{$f\sigma_8$ results: Gaussian vs.\ Student-$t$}

Table~\ref{tab:fsigma8_results} compares the two likelihoods on
analytical mocks. The Gaussian likelihood gives a $-71\%$ bias for
$S_\mathrm{P23}$; the Student-$t$ reduces this to $-2\%$.

\begin{table*}
\centering
\caption{$f\sigma_8$ results for the four scatter models (8 mocks
$\times$ 1500 SNe, $(f\sigma_8)_\mathrm{fid} = 0.428$).}
\label{tab:fsigma8_results}
\begin{tabular}{l|ccc|ccc}
\toprule
& \multicolumn{3}{c|}{Gaussian likelihood} &
\multicolumn{3}{c}{Student-$t$ likelihood} \\
Model &
$\langle f\sigma_8\rangle / \mathrm{fid}$ &
Error & Bias &
$\langle f\sigma_8\rangle / \mathrm{fid}$ &
Error & Bias \\
\midrule
$S_\mathrm{COH}$ & $0.65 \pm 0.35$ & 15\% & $-35\%$ & $0.88 \pm 0.14$ & 6\% & $-12\%$ \\
$S_\mathrm{G10}$ & $0.45 \pm 0.25$ & 13\% & $-55\%$ & $0.82 \pm 0.08$ & 3\% & $-18\%$ \\
$S_\mathrm{C11}$ & $0.40 \pm 0.29$ & 14\% & $-60\%$ & $0.82 \pm 0.10$ & 4\% & $-18\%$ \\
$S_\mathrm{P23}$ & $0.29 \pm 0.35$ & 15\% & $-71\%$ &
$\mathbf{0.98 \pm 0.27}$ & \textbf{12\%} & $\mathbf{-2\%}$ \\
\bottomrule
\end{tabular}
\end{table*}

\subsection{Results on 6600-SN Uchuu N-body mocks}

Table~\ref{tab:uchuu_results} shows the definitive test: 8 mocks
with 6600 SNe each on the Uchuu N-body velocity field. The mean
$f\sigma_8 = 0.462$ (fiducial 0.428), giving a bias of $+7.9\%$
with a mock-to-mock scatter of only 4.5\%. No mocks hit the
$f\sigma_8$ boundary. The coverage of 68\% CIs is 75\% (6/8).

\begin{table}
\centering
\caption{Uchuu N-body results with $S_\mathrm{P23}$ (8 mocks
$\times$ 6600 SNe, $(f\sigma_8)_\mathrm{fid} = 0.428$).}
\label{tab:uchuu_results}
\begin{tabular}{lccc}
\toprule
Seed & $f\sigma_8$ & 68\% CI & Covers fid. \\
\midrule
2000 & 0.464 & $[0.403, 0.535]$ & Y \\
2137 & 0.541 & $[0.481, 0.609]$ & N \\
2274 & 0.480 & $[0.421, 0.548]$ & Y \\
2411 & 0.417 & $[0.360, 0.482]$ & Y \\
2548 & 0.454 & $[0.388, 0.530]$ & Y \\
2685 & 0.453 & $[0.398, 0.515]$ & Y \\
2822 & 0.498 & $[0.436, 0.570]$ & N \\
2959 & 0.385 & $[0.318, 0.467]$ & Y \\
\midrule
Mean & 0.462 & & 6/8 \\
\bottomrule
\end{tabular}
\end{table}

\subsection{Comparison with Carreres et al.\ (2025)}

Table~\ref{tab:comparison} compares our results directly with
Table~4 of Carreres et al.

\begin{table}
\centering
\caption{Comparison with Carreres et al.\ (2025) for
$S_\mathrm{P23}$.}
\label{tab:comparison}
\begin{tabular}{lcc}
\toprule
Method & $\langle f\sigma_8\rangle / \mathrm{fid}$ & Bias \\
\midrule
True vel.\ fit (C25) & $0.988 \pm 0.013$ & $\sim 0\%$ \\
Simple fit, Gaussian (C25) & $0.808 \pm 0.036$ & $-20\%$ \\
BBC fit, Gaussian (C25) & $0.743 \pm 0.022$ & $-26\%$ \\
\midrule
\textbf{Student-$t$ (this work)} & $\bm{1.08 \pm 0.04}$ &
$\bm{+8\%}$ \\
\bottomrule
\end{tabular}
\end{table}

% ====================================================================
\section{Discussion}
\label{sec:discussion}

\subsection{Why the Student-$t$ works}

The P23 non-Gaussianity enters through the Tripp standardization:
the single $\beta$ coefficient cannot capture the dual-component
color model ($\beta_\mathrm{int} c_\mathrm{int} + (R_V{+}1)
E_\mathrm{dust}$). High-dust SNe receive biased distance estimates,
creating heavy-tailed outliers in the velocity distribution.

Crucially, the true peculiar velocities from gravitational dynamics
\emph{are} Gaussian (verified on Uchuu). The non-Gaussianity is
entirely an artifact of the SN~Ia standardization. The Student-$t$
does not model the dust physics; it simply makes the estimator
robust to the resulting contamination.

\subsection{Why BBC corrections failed}

We attempted four BBC-type corrections: (i) self-calibrating bin
means, (ii) analytical P23 corrections, (iii) split-$\beta$ Tripp
refitting, and (iv) color-dependent uncertainty inflation. All
increased the bias or degraded constraints. The fundamental issue
is that BBC corrects the \emph{mean} bias per bin, while the
$f\sigma_8$ bias arises from the \emph{shape} of the residual
distribution. Additionally, corrections risk removing velocity
signal through color--redshift correlations.

\subsection{Residual $+8\%$ bias on Uchuu}

The $+8\%$ positive bias on Uchuu mocks (vs.\ $-2\%$ on analytical
mocks) arises from the mismatch between the analytical velocity
power spectrum (Eisenstein-Hu + Bel et al.\ 2019 corrections) and
the actual N-body velocity correlations, which differ by $\sim
15\%$ at $k = 0.05$--$0.1\;h/\mathrm{Mpc}$. This is a limitation
of the covariance model, not the Student-$t$ likelihood. Using a
more accurate power spectrum (e.g.\ from CAMB/CLASS or measured
from the simulation) would reduce this further.

% ====================================================================
\section{Conclusions}
\label{sec:conclusions}

Replacing the Gaussian likelihood with a multivariate Student-$t$
in the $f\sigma_8$ analysis eliminates the P23 bias:

\begin{enumerate}
\item The P23 bias is reduced from $-20\%$ (Carreres et al.) to
$+8\%$ on Uchuu N-body mocks with 6600 SNe, and to $-2\%$ on
analytical mocks.

\item The method is fully data-driven: $\nu$ is estimated from the
kurtosis, requiring no knowledge of the scatter model.

\item The only change is one line in the log-likelihood
(Eq.~\ref{eq:nll_studentt} vs.\ Eq.~\ref{eq:gaussian}).

\item The method works on realistic N-body velocity fields, not
just analytical mocks.
\end{enumerate}

As SN~Ia samples grow to $\mathcal{O}(10^5)$ with LSST, the
Student-$t$ likelihood provides a practical, immediate solution for
peculiar velocity analyses with dust-based scatter models.

\section*{Acknowledgements}
This work used the Uchuu simulation (Ishiyama et al.\ 2021).

\end{document}
